% Texign book example: front matter, parts, chapters, appendix, back matter.
% Build with: texign build book.tex
% Defaults: sides=two, open=right, notes=foot (numbering restarts in every
% chapter), citations=academia.
\documentclass{texign-book}

\addbibresource{surse.bib}
\addterms{termeni.bib}

\subject{Exemplu Texign · Carte}
\title{Școala românească între lege și practică}
\subtitle{Surse, metode și interpretări (1864–1918)}
\author{István}
\publishers{București, 2026}
\date{noiembrie 2026}
\dedication{Învățătorilor de la sate, care au făcut școala posibilă.}

\newcommand\Para{Legea din 1864 a introdus învățământul primar obligatoriu și gratuit, dar aplicarea ei a depins de resursele fiecărei comune. Rapoartele revizorilor școlari arată o realitate inegală: în unele sate școala funcționa într-o casă închiriată, în altele lipsea cu totul.}
\newcommand\ParaB{Generația care a urmat a transformat școala într-un instrument al construcției naționale. Programele, manualele și formarea învățătorilor au devenit subiecte de dezbatere publică, iar presa a urmărit îndeaproape discuțiile parlamentare.}

\begin{document}
\frontmatter
\maketitle
\tableofcontents

\chapter{Prefață}
Această carte este un exemplu al clasei \texttt{texign-book}: arată paginile de titlu, părțile, capitolele, anexele, notele renumerotate la fiecare capitol, bibliografia, glosarul și indexul, toate în aceeași identitate vizuală ca articolele. Textul este demonstrativ; unele surse și cifre sunt fictive.

\mainmatter
\part{Surse}

\chapter{Legea și sursele ei}\label{cap:lege}

\Para\cite[art.~31]{lege1864} \ParaB\index{Legea instrucțiunii publice (1864)}

\section{Textul normativ}
\begin{source}{Legea instrucțiunii publice, 1864, art.~31}
„Instrucțiunea primară este obligatorie și gratuită.”
\end{source}

\Para\cite[45-48]{iorga1928}\index{Iorga, Nicolae} Legea creează o obligație fără să creeze și mijloacele ei.\note{Distincția dintre norma juridică și practica socială este un loc comun al istoriografiei educației.}

\begin{definition}[\Term{sursaprimara}]
Document produs în perioada studiată, de un participant sau de un martor al evenimentelor.
\end{definition}

\section{Rapoartele revizorilor}
Rapoartele \terms{revizor} sunt sursa cea mai apropiată de realitatea sălii de clasă.\cite{raport1872}\index{revizori școlari} \ParaB\cite[14]{popescu1929}

\image{revizor-1872.jpg}{Prima pagină a unui raport de revizor școlar, 1872 (ilustrație).}

\Para\cite[16]{popescu1929} \ParaB

\begin{table}[htbp]
\caption{Școli primare rurale și elevi înscriși, 1864–1914}
\begin{tabular*}{\linewidth}{@{\extracolsep{\fill}}lrrr@{}}
\toprule
\head{An} & \head{Școli} & \head{Elevi} & \head{Elevi / școală}\\
\midrule
1864 & 2\,036 & 61\,000 & 30\\
1890 & 2\,654 & 124\,000 & 47\\
1914 & 4\,730 & 553\,000 & 117\\
\bottomrule
\end{tabular*}
\end{table}

\chapter{Reformele lui Spiru Haret}\label{cap:haret}

\ParaB\cite{universul1906} Spiru Haret\index{Haret, Spiru} a legat explicit educația de dezvoltarea satului.\note{Spiru Haret (1851–1912), ministru al instrucțiunii în trei rânduri între 1897 și 1910.}

\begin{timeline}[Reformele școlare]
\event{1864}{Legea instrucțiunii publice}{Învățământul primar devine obligatoriu și gratuit.}
\event{1896}{Legea învățământului primar}{Obligativitate extinsă; \terms{scoalanormala} pentru învățători.}
\event{1897–1910}{Ministeriatele lui Spiru Haret}{Cercuri culturale, bănci populare, biblioteci sătești.}
\end{timeline}

\section{Școala și satul}
\Para\cite[112]{murgescu2010}\index{Murgescu, Bogdan} \ParaB\cite[114]{murgescu2010}

\begin{quotation}[Nicolae Iorga, \emph{Istoria învățământului românesc}, 1928]
Școala satului nu poate fi mai bună decât învățătorul ei.
\end{quotation}

\begin{comparison}[Legile din 1864 și 1896]{Criteriu}{1864}{1896}
\compare{Obligativitate}{8–12 ani}{7–14 ani}
\compare{Învățători}{fără formare obligatorie}{\terms{scoalanormala}}
\end{comparison}

\part{Interpretări}

\chapter{Limitele reformei}
\Para\cite[115]{murgescu2010} Răscoala din 1907\index{răscoala din 1907} a arătat cât de departe era satul de proiectul elitelor.\cite[Cf.][18-21]{bourdieu1970} \ParaB\cite[cap.~5]{hitchins1994}\index{Hitchins, Keith}

\begin{summary}
\begin{itemize}
\item Legea creează obligația; resursele vin mult mai târziu.
\item Reformele lui Haret mută accentul de la principiu la organizare.
\end{itemize}
\end{summary}

\appendix
\chapter{Documente}
\begin{source}{Legea instrucțiunii publice, 1864, art.~31}
„Instrucțiunea primară este obligatorie și gratuită. Părinții sau tutorii sunt datori a trimite copiii la școală de la vârsta de 8 până la 12 ani.”
\end{source}

\backmatter
\printsources
\printterms
\listoffigures
\listoftables
\printindex
\end{document}
